\documentclass[11pt,a4paper]{article}
\usepackage[T1]{fontenc}
\usepackage{isabelle,isabellesym}

\usepackage{amsmath}
\usepackage{amssymb}
\usepackage{stmaryrd}

% this should be the last package used
\usepackage{pdfsetup}

\urlstyle{rm}
\isabellestyle{it}

\begin{document}

\title{Operations on Laws of Continuous Paths}
\author{Dmitriy Traytel}
\maketitle

\begin{abstract}
  Arguments about controlled martingales on a space of continuous paths
  repeatedly cut a path at a time, continue it by another, or split it at a
  stopping time, and they need each of these operations not only on paths
  but on laws of paths.  This entry provides them for paths that carry a
  state in $\mathbb{R}^n$ together with its running covariation, on top of
  the path space $C([0,T])$ of the Archive entry
  \texttt{Continuous\_Path\_Spaces}.  Cutting, glueing, shifting, delaying
  and re-basing a path are continuous and Borel measurable for the
  supremum metric.  On laws, we construct the law of a concrete process,
  the glueing of a law with a measurable kernel of continuations chosen by
  the endpoint, and the additive glueing that splits a law at a stopping
  time and reassembles it; stopping times are treated as path functionals,
  with the regular conditional distribution of the future given the
  stopped past and optional sampling on the path space.  Generic lemmas
  carry the martingale identities, and the closed clauses of a constraint
  on the covariation, through weak limits of laws.  Finally, we define the
  class of martingale laws started at a point whose covariation has its
  density in a given set $S$ of matrices, derive its elementary properties
  under norm or diagonal bounds on $S$ (Lipschitz covariation paths, square
  integrability, second-moment bounds), and define the value of the
  minimum-exit-time problem over it, the supremum over the laws of the
  essential infimum of the exit time from a set $K$, with $S$ and $K$ as
  parameters.
\end{abstract}

\tableofcontents

\parindent 0pt\parskip 0.5ex

\input{session}

\bibliographystyle{abbrv}
\bibliography{root}

\end{document}
